\begin{textsample}{POS Dim 1 – profiled\_gpt – Score 6.00 – profiled\_gpt/t000532\_gpt.txt}  \label{ex:f1_pos_004}
I think if you give students really big, knotty historical topics, like the suffragettes or the world wars, you ’ve got far \textbf{more} chance of actually catching their curiosity.
There are so many \textbf{different} angles for them to hook into – the politics, the personalities, the military strategy, the social \textbf{changes} – that almost everyone can find something that genuinely interests them.
It ’s richer, \textbf{more} three‑dimensional, and it invites you to ask questions and make connections.
My \textbf{own} experience of history at school was nothing like that.
It was all chopped up into these very narrow, self‑contained little slices, and we were marched through them one by one.
There wasn't that \textbf{sense} of a big picture you could wander around in and explore from \textbf{different} directions.
Because of that, it never really lit that spark of curiosity for me ; it just felt like a series of disconnected topics.
The \textbf{way} it was taught, I mostly ended up just copying things off the board, rather than actually thinking about them or analysing what they meant.
I ’ve never been particularly good at straight memorisation, and no one ever really taught us how to memorise effectively either – it was just assumed you could do it.
Looking back, that feels like a real shame, especially considering how much that education cost.

% matched lemmas: change, different, more, own, sense, way
\end{textsample}
